\section{Conclusion}
\label{sec:conclusion}

We have presented the Neural Residual ODE (NXRO) framework for ENSO forecasting, systematically investigating how to combine physics-based structure with machine learning flexibility for climate prediction under severe data limitations. Our experiments, spanning 43 model variants evaluated on rigorous out-of-sample forecasting, yield insights relevant to both the climate science and machine learning communities.

\subsection{Principal Findings}

The central result of this work is that simpler hybrid architectures outperform both the full physics-based XRO model and pure deep learning approaches. The residual neural model, which combines a seasonal linear operator with a three-layer MLP, achieves the best deterministic forecast skill with test RMSE of 0.567°C---a 6.3\% improvement over the XRO baseline. This model omits the polynomial nonlinearities that are central to XRO's physics-based formulation, suggesting that the rigid polynomial basis may capture training-period patterns that do not generalize to future data.

The success of the linear-only variant, which outranks XRO despite having no nonlinear terms, reinforces this conclusion. With only 23 years of training data, the regularization provided by architectural simplicity outweighs the representational benefits of polynomial nonlinearities. The XRO model's variable-by-variable closed-form fitting, while computationally convenient, does not find globally optimal coefficients for multi-step forecasting. Joint gradient-based optimization, even with identical model structure, yields better generalization.

The graph-based and attention-based models demonstrate that structured sparsity provides effective regularization without sacrificing expressiveness. By restricting interactions to follow physics-informed patterns---either through fixed graph connectivity derived from XRO or through learned attention with masking---these models achieve competitive deterministic performance and excel at probabilistic forecasting.

The transformer-based model's failure ranks among our most important findings. Despite transformers' remarkable success in natural language processing~\cite{vaswani2017attention} and computer vision~\cite{dosovitskiy2020image}, the architecture performs worse than the physics-based XRO baseline on our climate forecasting task. This result carries implications for practitioners: modern deep learning architectures do not automatically transfer to small-data scientific domains. The approximately 300 monthly observations in our training set are insufficient for transformers to learn the seasonal structure, teleconnection patterns, and amplitude-dependent dynamics that the physics-informed linear operator encodes architecturally.

For probabilistic forecasting, likelihood-based noise optimization (Stage 2) consistently outperforms both post-hoc fitting and climate model-derived noise. The graph-based model with Stage 2 noise achieves CRPS of 0.488, representing a 13\% improvement over XRO. Equally important, the calibration improves substantially: spread-to-RMSE ratios increase from 0.53 (XRO) to 0.65 (NXRO-Graph), providing more reliable uncertainty quantification for decision-making applications.

The failure of climate model ensemble spread (Stage 3) to calibrate neural forecast uncertainty reveals a fundamental mismatch. The structural biases and parameterization differences captured by observation-simulation differences are irrelevant for well-optimized neural models that have already minimized prediction error. Uncertainty quantification methods must be matched to the properties of the prediction model.

\subsection{Broader Implications}

Our results speak to ongoing debates about the role of physics in machine learning for science~\cite{karniadakis2021physics,willard2022integrating}. The hybrid approach---encoding well-understood physics as architectural structure while learning corrections from data~\cite{rackauckas2020universal,yin2021augmenting}---emerges as the optimal strategy for the small-data regime characteristic of climate science and many other scientific domains.

Physics provides sample-efficient inductive bias~\cite{bronstein2021geometric}. The seasonal linear operator captures the dominant ENSO dynamics with far fewer parameters than a neural network would require to learn equivalent structure from data. This efficiency is critical when training samples number in the hundreds rather than millions, a regime where geometric and physical priors become essential for generalization.

Machine learning provides flexibility where physics is uncertain. The MLP residual, graph convolution, or attention mechanism can capture dynamics that the polynomial basis misses without requiring the modeler to specify functional forms a priori. The learned components adapt to the specific dataset rather than imposing assumptions that may not hold.

End-to-end optimization aligns training with the ultimate objective. By optimizing for multi-step forecast skill rather than one-step prediction error, the gradient-based approach finds coefficient configurations that minimize error accumulation over extended integration periods. This differs fundamentally from traditional fitting approaches that optimize each component independently.

The tension between model capacity and generalization requires careful navigation. Our experiments reveal that models achieving the best training performance often exhibit the largest generalization gaps. The art of hybrid modeling lies in providing sufficient flexibility to capture real dynamics while constraining the hypothesis space enough to prevent overfitting.

\subsection{Recommendations}

Based on our comprehensive evaluation, we offer guidance for practitioners developing ENSO forecast systems.

For operational deterministic forecasting, the residual neural model (NXRO-Res) trained directly on observations represents the recommended default. Its combination of physics-informed linear structure and flexible neural residual achieves the best out-of-sample skill while remaining computationally efficient and straightforward to train. Pre-training on climate model simulations does not improve and may degrade performance for this architecture.

For probabilistic forecasting, the graph-based model (NXRO-Graph) with Stage 2 noise optimization provides the best ensemble forecasts. The structured sparsity of the graph convolution appears particularly beneficial for uncertainty quantification, yielding both lower CRPS and better calibration than the residual model.

For applications requiring interpretable coupling patterns, the linear-only model (NXRO-Linear) or physics-structured model with frozen RO terms (NXRO-FixRO) preserve the interpretability of traditional approaches while benefiting from gradient-based optimization of the linear operator.

For two-stage transfer learning from climate model simulations, the pure neural ODE or physics-structured models benefit most from pre-training. The residual model's flexible MLP component may learn simulation-specific biases that are difficult to unlearn during fine-tuning on observations.

Regardless of architecture choice, rigorous out-of-sample evaluation is essential for model selection. In-sample metrics systematically favor complex models that overfit, providing misleading guidance about operational performance. Any evaluation of forecast systems should employ temporal holdout spanning multiple ENSO cycles.
